% ----------------------------------------------------------
% Apêndices
% ----------------------------------------------------------

% ---
% Inicia os apêndices
% ---
\begin{apendicesenv}

% Imprime uma página indicando o início dos apêndices
\partapendices

% ----------------------------------------------------------
\chapter{Algoritmos Utilizados} \label{sec:algorithm_append}

Os algoritmos aqui apresentados estão disponíveis em uma versão Python 3.10 no GitHub (\url{https://github.com/mephessivolc/tese})

\section{Algoritmo de Força Bruta}

O algoritmo de Força Bruta apresentado no \autoref{alg:brute_force} é utilizado ao longo de todo o projeto referente a algoritmo classico de força bruta.


\begin{algorithm}
\caption{Algoritmo de Força Bruta}
    \label{alg:brute_force}
    \begin{algorithmic}[1]
    \Require matriz de distâncias \texttt{dist}
    \Ensure melhor custo e melhor caminho
    \State $n \gets \text{len}(\texttt{dist})$
    \State \texttt{vertices $\gets$ [1, \dots, n-1]}
    \State \texttt{best\_cost $\gets$ inf}
    \State \texttt{best\_path $\gets$ None}
    \ForAll{\texttt{perm} em \texttt{permutations(vertices)}}
        \State \texttt{path $\gets$ (0,) + perm + (0,)}
        \State \texttt{cost $\gets$ 0}
        \For{$i \gets 0$ até $\text{len(path)} - 2$}
            \State \texttt{cost $\gets$ cost + dist[path[i]][path[i+1]]}
        \EndFor
        \If{\texttt{cost < best\_cost}}
            \State \texttt{best\_cost $\gets$ cost}
            \State \texttt{best\_path $\gets$ path}
        \EndIf
    \EndFor
    
    \State \Return \texttt{best\_cost, best\_path}
    \end{algorithmic}
\end{algorithm}


\section{Algoritmo Quântico}

\subsection{Equivalencia Qubits com Qumodes}

O \autoref{alg:hamiltoniano_qumode_binario} apresenta a construção do Hamiltoniano utilizado para codificar o TSP em qmodes. O mapeamento entre arestas dirigidas e modos bosônicos é definido no \autoref{alg:indice_aresta}. Após a obtenção de um estado candidato, o \autoref{alg:indices_para_matriz} converte as ocupações dos modos em uma matriz binária de arestas, enquanto o \autoref{alg:matriz_para_caminho} verifica se essa matriz representa um ciclo Hamiltoniano válido.

\begin{algorithm}[htbp]
\caption{Construção do Hamiltoniano qumode-binário para o TSP}
\label{alg:hamiltoniano_qumode_binario}
\begin{algorithmic}[1]
\Require Matriz de distâncias $d_{ij}$, penalidades $\alpha,\beta,\gamma$, dimensão local $\mathrm{dim}$, parâmetro $\mathrm{incluir\_p}$
\Ensure Hamiltoniano final $H_{\mathrm{final}}$

\State $N \gets$ número de cidades
\State $M \gets N(N-1)$ \Comment{um qumode para cada aresta dirigida $i \to j$, com $i \neq j$}

\For{$k \gets 0$ até $M-1$}
    \State construir o operador de destruição $a_k$
    \State $n_k \gets a_k^\dagger a_k$
\EndFor

\State $C \gets 0$
\For{$i \gets 0$ até $N-1$}
    \For{$j \gets 0$ até $N-1$}
        \If{$i \neq j$}
            \State $k \gets \Call{IndiceAresta}{i,j,N}$
            \State $C \gets C + d_{ij}n_k$
        \EndIf
    \EndFor
\EndFor

\State $H_1 \gets 0$
\For{$k \gets 0$ até $M-1$}
    \State $H_1 \gets H_1 + \alpha\left[n_k(n_k-1)\right]^2$
\EndFor

\State $H_2 \gets 0$
\For{$i \gets 0$ até $N-1$}
    \State $S_{\mathrm{saida}} \gets \sum_{j \neq i} n_{\mathrm{idx}(i,j)}$
    \State $H_2 \gets H_2 + \beta(S_{\mathrm{saida}}-1)^2$
\EndFor

\State $H_3 \gets 0$
\For{$j \gets 0$ até $N-1$}
    \State $S_{\mathrm{entrada}} \gets \sum_{i \neq j} n_{\mathrm{idx}(i,j)}$
    \State $H_3 \gets H_3 + \gamma(S_{\mathrm{entrada}}-1)^2$
\EndFor

\State $H_{\mathrm{final}} \gets C + H_1 + H_2 + H_3$

\If{$\mathrm{incluir\_p}$}
    \State construir o termo cinético $P$
    \State $H_{\mathrm{final}} \gets H_{\mathrm{final}} + P$
\EndIf

\State \Return $H_{\mathrm{final}}$
\end{algorithmic}
\end{algorithm}

\begin{algorithm}[htbp]
\caption{Mapeamento de uma aresta dirigida para o índice do qumode}
\label{alg:indice_aresta}
\begin{algorithmic}[1]
\Require Cidades $i,j$ e número de cidades $N$
\Ensure Índice $k$ associado à aresta $i \to j$

\If{$i=j$}
    \State lançar erro
\EndIf

\State $k \gets 0$
\For{$u \gets 0$ até $N-1$}
    \For{$v \gets 0$ até $N-1$}
        \If{$u \neq v$}
            \If{$u=i$ e $v=j$}
                \State \Return $k$
            \EndIf
            \State $k \gets k+1$
        \EndIf
    \EndFor
\EndFor
\end{algorithmic}
\end{algorithm}

\begin{algorithm}[htbp]
\caption{Conversão dos índices de ocupação em matriz binária de arestas}
\label{alg:indices_para_matriz}
\begin{algorithmic}[1]
\Require Tupla de ocupações $\mathrm{indices}$ e número de cidades $N$
\Ensure Matriz binária $X$ ou $\emptyset$

\State criar matriz $X \in \{0,1\}^{N \times N}$ inicialmente nula
\State $k \gets 0$

\For{$i \gets 0$ até $N-1$}
    \For{$j \gets 0$ até $N-1$}
        \If{$i \neq j$}
            \State $\mathrm{occ} \gets \mathrm{indices}[k]$

            \If{$\mathrm{occ} \notin \{0,1\}$}
                \State \Return $\emptyset$
            \EndIf

            \State $X_{ij} \gets \mathrm{occ}$
            \State $k \gets k+1$
        \EndIf
    \EndFor
\EndFor

\State \Return $X$
\end{algorithmic}
\end{algorithm}

\begin{algorithm}[htbp]
\caption{Validação da matriz binária como ciclo Hamiltoniano}
\label{alg:matriz_para_caminho}
\begin{algorithmic}[1]
\Require Matriz binária de arestas $X$
\Ensure Caminho fechado válido ou $\emptyset$

\State $N \gets$ dimensão de $X$

\If{alguma linha de $X$ não soma $1$}
    \State \Return $\emptyset$
\EndIf

\If{alguma coluna de $X$ não soma $1$}
    \State \Return $\emptyset$
\EndIf

\State $\mathrm{caminho} \gets [0]$
\State $\mathrm{atual} \gets 0$
\State $\mathrm{visitados} \gets \{0\}$

\For{$t \gets 1$ até $N-1$}
    \State encontrar $\mathrm{prox}$ tal que $X_{\mathrm{atual},\mathrm{prox}}=1$

    \If{$\mathrm{prox}$ não existe ou não é único}
        \State \Return $\emptyset$
    \EndIf

    \If{$\mathrm{prox} \in \mathrm{visitados}$}
        \State \Return $\emptyset$
    \EndIf

    \State adicionar $\mathrm{prox}$ ao caminho
    \State adicionar $\mathrm{prox}$ a $\mathrm{visitados}$
    \State $\mathrm{atual} \gets \mathrm{prox}$
\EndFor

\State encontrar $\mathrm{prox}$ tal que $X_{\mathrm{atual},\mathrm{prox}}=1$

\If{$\mathrm{prox} \neq 0$}
    \State \Return $\emptyset$
\EndIf

\State adicionar $0$ ao final de $\mathrm{caminho}$

\If{$|\mathrm{visitados}| \neq N$}
    \State \Return $\emptyset$
\EndIf

\State \Return $\mathrm{caminho}$
\end{algorithmic}
\end{algorithm}

\subsection{Estados de Fock truncado na posição de visita de cada cidade}

Os \autoref{alg:projector_k_a}---\autoref{alg:analisar_estados_base} descrevem a construção e a análise da formulação do TSP em qumodes truncados. O \autoref{alg:projector_k_a} define os projetores locais $P_{k,a}$, que identificam a ocupação do nível $a$ no registrador associado à posição $k$ do caminho. Esses projetores são empregados no \autoref{alg:hamiltoniano_tsp_qudit} para construir o Hamiltoniano do problema, composto por um termo de distância, $H_{\mathrm{dist}}$, responsável por codificar o custo do percurso, e por um termo de penalização, $H_{\mathrm{pen}}$, que evita a repetição de cidades. Em seguida, o \autoref{alg:validacao_caminho} verifica se uma sequência de índices constitui uma permutação válida das cidades e a converte em um ciclo fechado. Finalmente, o \autoref{alg:analisar_estados_base} percorre os estados da base computacional, calcula suas energias esperadas e identifica a configuração válida de menor energia.

\begin{algorithm}[htbp]
\caption{Construção do projetor local}
\label{alg:projector_k_a}
\begin{algorithmic}[1]
\Require posição $k$, nível $a$, número de registradores $N$, dimensão local $\mathrm{dim}$
\Ensure operador projetor $P_{k,a}$

\State $\mathrm{ops} \gets [\ ]$
\For{$m \gets 0$ até $N-1$}
    \If{$m = k$}
        \State adicionar $|a\rangle\langle a|$ em $\mathrm{ops}$
    \Else
        \State adicionar $\mathbb{I}_{\mathrm{dim}}$ em $\mathrm{ops}$
    \EndIf
\EndFor
\State \Return produto tensorial dos elementos de $\mathrm{ops}$
\end{algorithmic}
\end{algorithm}

\begin{algorithm}[htbp]
\caption{Construção do Hamiltoniano do TSP em qumodes}
\label{alg:hamiltoniano_tsp_qudit}
\begin{algorithmic}[1]
\Require matriz de distâncias $d_{ij}$, penalização $\lambda$, dimensão local $\mathrm{dim}$
\Ensure Hamiltoniano $H$

\State $N \gets$ número de linhas de $d_{ij}$

\If{$\mathrm{dim}$ não informado}
    \State $\mathrm{dim} \gets N$
\EndIf

\If{$\mathrm{dim} < N$}
    \State lançar erro
\EndIf

\For{$k \gets 0$ até $N-1$}
    \For{$a \gets 0$ até $N-1$}
        \State $P_{k,a} \gets \Call{ProjetorLocal}{k,a,N,\mathrm{dim}}$
    \EndFor
\EndFor

\State $H_{\mathrm{dist}} \gets 0$
\For{$k \gets 0$ até $N-1$}
    \State $k' \gets (k+1)\bmod N$
    \For{$a \gets 0$ até $N-1$}
        \For{$b \gets 0$ até $N-1$}
            \State $H_{\mathrm{dist}} \gets H_{\mathrm{dist}} + d_{ij}[a,b]\,P_{k,a}P_{k',b}$
        \EndFor
    \EndFor
\EndFor

\State $H_{\mathrm{pen}} \gets 0$
\For{$k \gets 0$ até $N-1$}
    \For{$l \gets k+1$ até $N-1$}
        \For{$a \gets 0$ até $N-1$}
            \State $H_{\mathrm{pen}} \gets H_{\mathrm{pen}} + P_{k,a}P_{l,a}$
        \EndFor
    \EndFor
\EndFor

\State \Return $H_{\mathrm{dist}} + \lambda H_{\mathrm{pen}}$
\end{algorithmic}
\end{algorithm}

\begin{algorithm}[htbp]
\caption{Validação e conversão de índices em caminho}
\label{alg:validacao_caminho}
\begin{algorithmic}[1]
\Require sequência $\mathrm{indices}$, número de cidades $N$
\Ensure caminho válido ou valor nulo

\If{$|\mathrm{indices}| \neq N$}
    \State \Return falso
\EndIf

\If{$\mathrm{sorted}(\mathrm{indices}) \neq [0,1,\dots,N-1]$}
    \State \Return falso
\EndIf

\State \Return tupla $(\mathrm{indices}[0], \ldots, \mathrm{indices}[N-1], \mathrm{indices}[0])$
\end{algorithmic}
\end{algorithm}

\begin{algorithm}[htbp]
\caption{Análise dos estados da base computacional}
\label{alg:analisar_estados_base}
\begin{algorithmic}[1]
\Require Hamiltoniano $H$, matriz de distâncias $\mathrm{dist}$, número de cidades $N$, dimensão local $\mathrm{dim}$
\Ensure melhor energia, melhor caminho e lista de resultados válidos

\If{$\mathrm{dim}$ não informado}
    \State $\mathrm{dim} \gets N$
\EndIf

\State $\mathrm{melhor\_energia} \gets +\infty$
\State $\mathrm{melhor\_indices} \gets \emptyset$
\State $\mathrm{melhor\_path} \gets \emptyset$
\State $\mathrm{resultados\_validos} \gets [\ ]$

\ForAll{$\mathrm{indices} \in \{0,\dots,\mathrm{dim}-1\}^N$}
    \If{$\mathrm{indices}$ não formam uma permutação válida de $\{0,\dots,N-1\}$}
        \State \textbf{continue}
    \EndIf

    \State construir o estado tensorial correspondente
    \State calcular a energia esperada de $H$
    \State converter $\mathrm{indices}$ em caminho fechado
    \State calcular o custo do caminho
    \State armazenar resultado em $\mathrm{resultados\_validos}$

    \If{$\mathrm{energia} < \mathrm{melhor\_energia}$}
        \State atualizar melhor energia
        \State atualizar melhores índices
        \State atualizar melhor caminho
    \EndIf
\EndFor

\State \Return melhores resultados
\end{algorithmic}
\end{algorithm}




\end{apendicesenv}
% ---